% DESIGN29 working paper. Every measured number is a macro from numbers.tex,
% generated by make_numbers.py from research/*.json; `make check` fails if
% numbers.tex is stale or any reference does not resolve.
\documentclass[11pt]{article}
\usepackage[margin=1.1in]{geometry}
\usepackage{amsmath,amssymb}
\usepackage{booktabs}
\usepackage{microtype}
\usepackage[round]{natbib}
\usepackage[hidelinks]{hyperref}
\setlength{\abovecaptionskip}{4pt}
\setlength{\belowcaptionskip}{8pt}
% Generated by make_numbers.py from research/*.json. Do not edit by hand.
\newcommand{\mmMktA}{4.93}
\newcommand{\mmMktSE}{1.57}
\newcommand{\mmMktB}{0.60}
\newcommand{\mmMktN}{1,065}
\newcommand{\mmMomA}{12.83}
\newcommand{\mmMomSE}{1.72}
\newcommand{\mmMomB}{0.45}
\newcommand{\mmMomN}{1,060}
\newcommand{\mmSMBA}{\ensuremath{-0.59}}
\newcommand{\mmSMBSE}{0.91}
\newcommand{\mmSMBB}{0.60}
\newcommand{\mmSMBN}{1,065}
\newcommand{\mmHMLA}{1.73}
\newcommand{\mmHMLSE}{1.05}
\newcommand{\mmHMLB}{0.54}
\newcommand{\mmHMLN}{1,065}
\newcommand{\mmRMWA}{2.80}
\newcommand{\mmRMWSE}{0.89}
\newcommand{\mmRMWB}{0.58}
\newcommand{\mmRMWN}{621}
\newcommand{\mmCMAA}{0.36}
\newcommand{\mmCMASE}{0.67}
\newcommand{\mmCMAB}{0.67}
\newcommand{\mmCMAN}{621}
\newcommand{\tabMM}{Mkt & 4.86 & (1.56) & 0.61 & 4.93 & (1.57) & 0.60 & 1,065 & 1.17 \\
UMD & 12.51 & (1.71) & 0.47 & 12.83 & (1.72) & 0.45 & 1,060 & 1.28 \\
SMB & \ensuremath{-0.58} & (0.91) & 0.62 & \ensuremath{-0.59} & (0.91) & 0.60 & 1,065 & 0.91 \\
HML & 1.97 & (1.02) & 0.57 & 1.73 & (1.05) & 0.54 & 1,065 & 1.02 \\
RMW & 2.44 & (0.83) & 0.62 & 2.80 & (0.89) & 0.58 & 621 & 0.83 \\
CMA & 0.38 & (0.67) & 0.68 & 0.36 & (0.67) & 0.67 & 621 & 0.67 \\}
\newcommand{\cowyMktSRo}{0.42}
\newcommand{\cowyMktSRm}{0.52}
\newcommand{\cowyMktD}{\ensuremath{+0.10}}
\newcommand{\cowyMktComb}{0.44}
\newcommand{\cowyMktOrig}{0.46}
\newcommand{\cowyMktTD}{\ensuremath{-0.02}}
\newcommand{\cowyMktP}{0.80}
\newcommand{\cowyMomSRo}{0.48}
\newcommand{\cowyMomSRm}{0.99}
\newcommand{\cowyMomD}{\ensuremath{+0.51}}
\newcommand{\cowyMomComb}{0.93}
\newcommand{\cowyMomOrig}{0.48}
\newcommand{\cowyMomTD}{\ensuremath{+0.45}}
\newcommand{\cowyMomP}{0.00}
\newcommand{\cowyMktSpan}{1936-08--2016-12}
\newcommand{\ltzDDrt}{\ensuremath{-90.8}}
\newcommand{\ltzDDfs}{\ensuremath{-54.4}}
\newcommand{\ltzDDmkt}{\ensuremath{-50.3}}
\newcommand{\ltzA}{3.40}
\newcommand{\ltzSE}{2.30}
\newcommand{\ltzT}{1.5}
\newcommand{\ltzLev}{18.9}
\newcommand{\ltzSpan}{1936-08--2016-12}
\newcommand{\mktUniA}{2.67}
\newcommand{\mktUniT}{2.94}
\newcommand{\mktUniB}{0.39}
\newcommand{\mktUniSpan}{1927-07--2012-12}
\newcommand{\mktUniN}{1,026}
\newcommand{\mktFullA}{\ensuremath{-1.19}}
\newcommand{\mktFullT}{\ensuremath{-1.18}}
\newcommand{\mktFullRsq}{0.73}
\newcommand{\mktFullSpan}{1963-07--2012-12}
\newcommand{\mktFullBMkt}{0.62}
\newcommand{\mktFullBSMB}{0.01}
\newcommand{\mktFullBHML}{\ensuremath{-0.07}}
\newcommand{\mktFullBRMW}{0.10}
\newcommand{\mktFullBCMA}{0.06}
\newcommand{\mktFullBMom}{0.10}
\newcommand{\mktFullBBAB}{0.06}
\newcommand{\mktSRm}{0.53}
\newcommand{\mktSRr}{0.39}
\newcommand{\mktSRd}{\ensuremath{+0.14}}
\newcommand{\mktLWp}{0.047}
\newcommand{\mktRWp}{0.86}
\newcommand{\rwSpan}{1965-07--2012-12}
\newcommand{\rwN}{570}
\newcommand{\spaP}{0.0002}
\newcommand{\tabEras}{1927--1945 & 4.31 & (1.97) & 222 \\
1946--1962 & 1.59 & (1.85) & 204 \\
1963--1987 & 0.76 & (0.64) & 300 \\
1988--2012 & 1.04 & (0.79) & 300 \\}
\newcommand{\erasPos}{4}
\newcommand{\swStat}{8.67}
\newcommand{\swAt}{1998-04}
\newcommand{\swP}{0.051}
\newcommand{\verdict}{NO ALPHA SURVIVES}
\newcommand{\critOne}{fails}
\newcommand{\critTwo}{holds}
\newcommand{\critThree}{fails}
\newcommand{\critFour}{holds}
\newcommand{\turnover}{5.25}
\newcommand{\breakevenUni}{51}
\newcommand{\breakevenFull}{\ensuremath{-23}}
\newcommand{\meanExposure}{0.71}
\newcommand{\shareCap}{0.8}
\newcommand{\emHalfTickBp}{0.21}
\newcommand{\gridSRmin}{0.457}
\newcommand{\gridSRmax}{0.530}
\newcommand{\gridSRprim}{0.530}
\newcommand{\gridPBO}{0.80}
\newcommand{\gridCombos}{12,870}
\newcommand{\dsrProg}{1.00}
\newcommand{\dsrSRzero}{0.106}
\newcommand{\nTrials}{202}
\newcommand{\repoPBO}{0.45}
\newcommand{\repoDSR}{0.82}
\newcommand{\repoVariants}{140}
\newcommand{\dtbUniA}{1.19}
\newcommand{\dtbUniT}{1.35}
\newcommand{\dtbFullA}{\ensuremath{-1.23}}
\newcommand{\dtbFullT}{\ensuremath{-1.22}}
\newcommand{\dtbSpread}{0.35}
\newcommand{\conA}{\ensuremath{-0.26}}
\newcommand{\conT}{\ensuremath{-0.2}}
\newcommand{\conSRm}{0.73}
\newcommand{\conSRr}{0.91}
\newcommand{\conSpan}{2013-01--2026-08}
\newcommand{\vixHold}{0.86}
\newcommand{\vixCon}{0.86}
\newcommand{\famSMBA}{0.99}
\newcommand{\famSMBT}{1.14}
\newcommand{\famHMLA}{5.78}
\newcommand{\famHMLT}{4.00}
\newcommand{\famRMWA}{7.43}
\newcommand{\famRMWT}{4.66}
\newcommand{\famCMAA}{4.80}
\newcommand{\famCMAT}{3.54}
\newcommand{\famMomA}{8.81}
\newcommand{\famMomT}{3.97}
\newcommand{\famBABA}{12.18}
\newcommand{\famBABT}{5.53}
\newcommand{\tabFamily}{SMB & 0.68 & (1.03) & 0.99 & (1.14) & 0.209 & 0.15 & 0.21 & 0.199 & 1.6 \\
HML & 4.06 & (3.68) & 5.78 & (4.00) & $<$0.001 & 0.43 & 0.61 & 0.014 & 2.0 \\
RMW & 5.50 & (3.53) & 7.43 & (4.66) & $<$0.001 & 0.48 & 0.75 & 0.015 & 3.3 \\
CMA & 2.91 & (2.82) & 4.80 & (3.54) & $<$0.001 & 0.60 & 0.71 & 0.083 & 2.7 \\
UMD & 7.88 & (5.85) & 8.81 & (3.97) & $<$0.001 & 0.46 & 0.88 & $<$0.001 & 1.6 \\
BAB & 8.57 & (5.33) & 12.18 & (5.53) & $<$0.001 & 0.81 & 1.19 & $<$0.001 & 1.9 \\}
\newcommand{\scaleMin}{1.6}
\newcommand{\scaleMax}{3.3}


\title{Is a Production Volatility-Targeting Rule an Alpha?\\
A Sealed, Pre-Registered Test on a Century of Factor Returns}
\author{}
\date{Working paper --- October 2026}

\begin{document}
\maketitle

\begin{abstract}
\noindent
Volatility-managed portfolios earn large spanning alphas
\citep{moreira2017}. Those alphas are hard to capture in real time
\citep{cederburg2020}, and part of them reflects a leverage constant
estimated with look-ahead \citep{liu2019}. We ask a narrower question with a
stricter design: does a specific, already-deployed volatility-targeting
overlay earn alpha out of the era it was designed in? The overlay combines a
HAR volatility-forecast lever capped at $1.5\times$, CVaR and drawdown
brakes, and a momentum tilt, and it is causal by construction. We
pre-registered the test, data, tolerances and decision rule; first
replicated the published results within stated tolerances; then sealed the
1926--2012 holdout by committing a hash of the code before running it once.

On the market, the overlay raises the Sharpe ratio from \mktSRr{} to
\mktSRm{} (Ledoit--Wolf one-sided $p=\mktLWp$), and its univariate spanning
alpha is \mktUniA\% a year ($t=\mktUniT$). Net of the Fama--French five
factors, momentum and betting-against-beta, the alpha is \mktFullA\%
($t=\mktFullT$, \mktFullSpan). The managed market loads on profitability,
momentum and BAB, and the pre-registered verdict is that no alpha survives.

Applied to factor portfolios, the same real-time rule earns gross
full-model alphas that survive Romano--Wolf correction for HML, RMW, CMA,
momentum and BAB. These are gross of the factors' own trading costs, which
factor returns cannot price. We state what stock-level CRSP and TAQ would
change and stop there.

\medskip
\noindent\textbf{JEL:} G11, G12, C12, C58.\quad
\textbf{Keywords:} volatility targeting, volatility-managed portfolios,
pre-registration, spanning tests, multiple testing, backtest overfitting.
\end{abstract}

\section{Introduction}

A volatility-targeting rule scales exposure to a risky portfolio inversely to
a forecast of its volatility. \citet{fleming2001} document the economic value
of volatility timing in asset allocation. \citet{barroso2015} show that
scaling momentum by its own realized variance roughly doubles its Sharpe
ratio and removes its crashes \citep[see also][]{daniel2016}.
\citet{harvey2018} find that volatility targeting raises Sharpe ratios for
equities and reduces the probability of extreme returns.

\citet{moreira2017} make the strongest form of the claim. They scale a
factor by $c/\hat\sigma^2_{t-1}$, where $\hat\sigma^2_{t-1}$ is last month's
realized variance, and regress the result on the unmanaged factor. The
spanning alpha is large for the market and for most factors.

The literature since has pushed back on whether investors can collect it.
\citet{cederburg2020} build the strategies an investor could have run in
real time. They find that combining the managed and original factor improves
on the original for some factors, momentum and BAB above all, but not for
the market. \citet{liu2019} point out that the constant $c$ is chosen on the
full sample. With $c$ estimated in real time, they report maximum drawdowns
that are deeper than the market's.

We study one deployed rule rather than the textbook one. It is the risk
overlay of an open-source quantitative platform, and it was designed and
tuned on 2013--2026 data. It differs from the \citet{moreira2017} rule in
ways a practitioner would recognize:
\begin{itemize}
\item it targets an annualized volatility of 13\% using a causal HAR
  forecast \citep{corsi2009}, refit monthly;
\item its leverage is capped at $1.5\times$;
\item CVaR and drawdown brakes cut exposure after losses;
\item the result is smoothed over five days;
\item a $\pm15\%$ tilt follows the sign of the trailing 21-day return.
\end{itemize}
Every input is trailing, so the rule needs no full-sample constant.

The question is whether this rule earns alpha, net of the factors a referee
would name, on years it has never seen.

Three design choices make the answer credible.
\begin{enumerate}
\item \emph{Pre-registration and a seal.} Before any file was downloaded, we
  committed:
  \begin{itemize}
  \item the question and the data;
  \item replication targets with numeric tolerances;
  \item the test battery;
  \item a four-part decision rule.
  \end{itemize}
  The 1926--2012 holdout was then sealed. The SHA-256 of the exact code that
  would produce the verdict was committed before that code first touched a
  pre-2013 return, and the code was run once.
\item \emph{Replicate first.} The program was allowed to proceed only if the
  published results of \citet{moreira2017} and \citet{cederburg2020}
  reproduced within the registered tolerances on today's data. They did.
\item \emph{Survivorship-free test assets.} The overlay is a portfolio-level
  timing rule, so we test it on portfolios rather than on a hand-picked
  stock universe. We use Ken French's market and factor portfolios, built
  from CRSP from 1926, and AQR's BAB factor.
\end{enumerate}

\paragraph{Findings.}
\begin{itemize}
\item \emph{The market.} On the market, the overlay is a real Sharpe-ratio
  improvement for an investor who holds only the market: \mktSRm{} against
  \mktSRr{} over \mktUniSpan. It is not an alpha. Net of FF5, momentum and
  BAB the alpha is \mktFullA\% a year ($t=\mktFullT$). The managed market's
  extra return is explained by its loadings on profitability (RMW), momentum
  and BAB. Most of the univariate alpha accrues before 1963.
\item \emph{Costs.} Trading costs do not decide this result. The breakeven
  cost on the univariate alpha is \breakevenUni{}~bp per unit of turnover,
  against about \emHalfTickBp{}~bp for half an E-mini tick today.
\item \emph{Factor portfolios.} On factor portfolios the same rule earns
  large gross alphas in real time, and they survive Romano--Wolf
  correction. We do not claim them net of costs.
\end{itemize}

\section{Data}

\paragraph{Sources.}
\begin{itemize}
\item \emph{Factors.} Daily and monthly Fama--French research factors and
  momentum from Ken French's data library, built from the 202608 CRSP
  database: Mkt$-$RF, SMB, HML and RF from 1926-07; RMW and CMA from
  1963-07 \citep{fama2015}; UMD from 1926-11 \citep{carhart1997}.
\item \emph{BAB.} AQR's daily US betting-against-beta factor, from 1930-12
  \citep{frazzini2014}.
\item \emph{VIX.} Cboe's VIX, used only descriptively.
\end{itemize}

\paragraph{Archiving.} French revises history, so a URL does not identify a
dataset. Every file was fetched once, after the registration was committed,
and is pinned by SHA-256 in a committed manifest. Every loader verifies the
hash before parsing and refuses a mismatch. The raw files are not
redistributed. Anyone reproducing either obtains the same bytes or is told
that they hold a different vintage.

CSV floats are parsed with correctly rounded (round-trip) conversion. The
default parser in our stack rounds through long double, which differs
between x86 and arm64 hardware.

\paragraph{Financing.} The factors are excess returns over French's
one-month T-bill rate RF. Holding exposure $e$ in Mkt$-$RF is the same as
holding $e$ in the market and $1-e$ in bills, or borrowing $e-1$ at the bill
rate. Leverage is therefore financed at the T-bill rate. A robustness check
finances it at the three-month bill (DTB3) instead.

\paragraph{Monthly returns.} The overlay rebalances daily. The monthly excess
return of a position with daily exposure $e_t$ is
\[
R = \prod_t (1 + e_{t-1} f_t + r^f_t) - \prod_t (1 + r^f_t),
\]
taken over a calendar month's trading days. Every managed and unmanaged
series in the Stage-2 tests is built this way, so no comparison mixes
French's monthly files with daily-compounded returns.

\paragraph{Saturdays and holidays.} The NYSE traded on Saturdays until 1952,
and the overlay's windows assume weekday sessions. Each Saturday return is
therefore compounded into the next observation. AQR's calendar lists 382 US
holidays from 1984 on that French's calendar omits; BAB returns on those days
are compounded forward the same way. The replication uses French's files
unchanged, as the replicated papers did.

\section{Design}

\subsection{Windows and the seal}
\begin{itemize}
\item \emph{Holdout:} 1926-07 to 2012-12.
\item \emph{Full-model span:} 1963-07 to 2012-12, the span over which all
  seven regressors exist.
\item \emph{Contaminated window:} 2013-01 onward. The overlay was designed
  on these years, so they are reported descriptively and enter no decision.
\end{itemize}

The registration, an amendment fixing the replication samples, and the seal
were committed in that order, each before the computation it governs. The
seal amendment records the code digest. A verification gate refuses any
results file whose recorded digest differs.

\subsection{Stage 1: the replication gate}

\paragraph{Construction.} The managed factor is
$f^\sigma_t = (c/\mathrm{RV}^2_{t-1}) f_t$, where:
\begin{itemize}
\item $\mathrm{RV}^2$ is the demeaned sum of squared daily returns within
  the month;
\item $c$ equates the standard deviations of $f^\sigma$ and $f$ over the
  sample.
\end{itemize}

\paragraph{Targets and tolerances.} We register the Table~1 spanning alphas
of \citet{moreira2017}, with heteroskedasticity-robust standard errors. The
samples match their reported $N$ exactly: they end in 2015-04 and start in
1926-08, 1927-01 or 1963-08.

The tolerance on $|\Delta\alpha|$ is $0.75\,\mathrm{SE}$ for the market and
momentum, plus $|\Delta\beta|\le0.06$ for those two; it is
$1\,\mathrm{SE}$ for SMB, HML, RMW and CMA. The tolerances allow for French's
revisions since the paper was written.

From \citet{cederburg2020} we register:
\begin{itemize}
\item the in-sample Sharpe ratios (their Table~1);
\item the real-time combination strategy (their Table~5), built as follows:
  \begin{itemize}
  \item an expanding window that starts after $K=120$ months;
  \item $c_t$ estimated on the training window;
  \item mean--variance weights at $\gamma=5$;
  \item the position clipped at $|y_t|\le5$.
  \end{itemize}
\end{itemize}

\paragraph{Pass rule.} The gate passes if:
\begin{itemize}
\item the market and momentum both pass;
\item at least three of the other four factors pass;
\item all three Cederburg et al.\ market checks pass.
\end{itemize}
If the gate failed, the program stopped.

\paragraph{The look-ahead point.} The full text of \citet{liu2019} was not
available to us. We register only their abstract's claim, as a descriptive
check: maximum drawdowns of roughly 68--93\% once $c$ is real-time. We
reproduce the construction directly.

\subsection{Stage 2: the overlay}

\paragraph{Primary.} The primary strategy is the overlay on the market,
exactly as deployed:
\begin{itemize}
\item \emph{Lever:} target 13\%, HAR forecast, cap $1.5\times$.
\item \emph{CVaR brake:} target 1.8\% daily at 95\%.
\item \emph{Drawdown brake:} a linear ramp from $-8\%$ to $-20\%$, with a
  floor of 0.25.
\item \emph{Smoothing:} 5 days.
\item \emph{Tilt:} $\pm15\%$.
\end{itemize}
Its estimators read the market's total return, as they read the deployed
book's. Exposure decided at the close of day $t$ applies to day $t+1$. The
first 252 observations train the estimators and are not evaluated.

A verification gate shows that this code reproduces the deployed backtest's
applied exposure exactly, on every live day of a synthetic panel.

One property of the rule was known in advance and registered. The drawdown
brake reads the market's own high-water mark, so after 1929 it holds near
its floor until that mark is regained.

\paragraph{Family.} The family applies the same rule to SMB, HML, RMW, CMA,
UMD and BAB. The deployed thresholds are denominated in the deployed book's
volatility, so each factor is first expressed in those units:
$\tilde f_t = f_t \cdot 0.13/\hat\sigma_{t-1}$. Here $\hat\sigma$ is the
factor's trailing realized volatility over 2{,}520 observations, expanding
from a 252-observation minimum. No alternative normalization was tried.

\paragraph{Tests.} All regressions are monthly, with \citet{newey1987}
standard errors at 12 lags and alphas annualized. The battery is:
\begin{itemize}
\item univariate spanning regressions on each managed factor's own unmanaged
  version;
\item full-model regressions on Mkt, SMB, HML, RMW, CMA, UMD and BAB;
\item the \citet{ledoit2008} studentized circular-block bootstrap test of the
  Sharpe-ratio difference (12-month blocks, 5{,}000 resamples, one-sided);
\item the \citet{romano2005} studentized stepdown across the seven
  full-model alphas, at 5\% family-wise;
\item the SPA test of \citet{hansen2005} on the factor-hedged series, using the
  stationary bootstrap of \citet{politis1994};
\item univariate alphas in four eras;
\item an \citet{andrews1993} sup-Wald test for a break in the intercept,
  with its $p$-value simulated from Brownian bridges.
\end{itemize}

\paragraph{Overfitting diagnostics.} A registered eight-member grid of
deployed-rule siblings (tilt, cap and brakes, each on or off) supplies:
\begin{itemize}
\item a CSCV probability of backtest overfitting \citep{bailey2017};
\item the dispersion for a deflated Sharpe ratio \citep{bailey2014}.
\end{itemize}
Every member and every family factor is charged to the platform's
cumulative trial ledger, which now stands at $N=\nTrials$.

\paragraph{Decision rule.} The overlay earns alpha if and only if, for the
market, all four of the following hold:
\begin{enumerate}
\item the full-model alpha is positive with $t\ge3$, the hurdle of
  \citet{harvey2016};
\item the Ledoit--Wolf $p<0.05$;
\item the alpha survives Romano--Wolf;
\item the univariate alpha is positive in at least three of four eras.
\end{enumerate}

\section{Stage 1: the published results replicate}

Table~\ref{tab:mm} compares the replication with \citet{moreira2017}. Every
factor lands inside its tolerance, most by a wide margin. The largest gap
relative to its tolerance is RMW's: \mmRMWA{} against a published 2.44.

\begin{table}[t]
\centering\small
\caption{Moreira--Muir Table~1 on the 202608 French vintage. Spanning
regression $f^\sigma=\alpha+\beta f+\varepsilon$, monthly, $\alpha$ in \% a
year, White SEs in parentheses. Tolerance: the registered bound on
$|\Delta\alpha|$.}
\label{tab:mm}
\begin{tabular}{lrrrrrrrr}
\toprule
& \multicolumn{3}{c}{Published} & \multicolumn{4}{c}{Replication} & \\
\cmidrule(lr){2-4}\cmidrule(lr){5-8}
Factor & $\alpha$ & (SE) & $\beta$ & $\alpha$ & (SE) & $\beta$ & $N$ & Tol. \\
\midrule
\tabMM
\bottomrule
\end{tabular}
\end{table}

The \citet{cederburg2020} checks also pass.
\begin{itemize}
\item \emph{In sample.} The managed market's Sharpe ratio is \cowyMktSRm{}
  against \cowyMktSRo{}, as published (0.51 against 0.42). The momentum gain
  is \cowyMomD{} (published $+0.50$).
\item \emph{Real time, market.} The combination strategy over
  \cowyMktSpan{} has a Sharpe ratio of \cowyMktComb{}, against \cowyMktOrig{}
  for the original factor alone. The difference is \cowyMktTD{}, with
  $p=\cowyMktP$ in the test of \citet{jobson1981}. The published values are
  0.42, 0.46 and $-0.04$.
\item \emph{Real time, momentum.} The gain survives in real time: \cowyMomComb{}
  against \cowyMomOrig{}, as published (0.92 against 0.49).
\end{itemize}

The look-ahead point also reproduces. With $c$ estimated in real time
(\ltzSpan), the managed market's spanning alpha falls to \ltzA\%
(SE \ltzSE, $t=\ltzT$). Leverage reaches \ltzLev$\times$, and the maximum
drawdown is \ltzDDrt\%, against \ltzDDfs\% with the full-sample constant and
\ltzDDmkt\% for the market. That lies inside the 68--93\% range
\citet{liu2019} report.

The market alpha of a volatility-managed portfolio thus depends on the
full-sample constant. The rule we test has no such constant.

\section{Stage 2: the deployed overlay}

\subsection{The market}

The verdict is \textbf{\verdict}. Of the four registered criteria:
\begin{itemize}
\item the full-model $t\ge3$ test \critOne;
\item the Ledoit--Wolf test \critTwo;
\item Romano--Wolf survival \critThree;
\item positive alpha in at least three of four eras \critFour.
\end{itemize}

For an investor holding only the market, the overlay is an improvement.
Over \mktUniSpan{} ($N=\mktUniN$ months) its Sharpe ratio is \mktSRm{}
against \mktSRr{}, a gain of \mktSRd{} with Ledoit--Wolf one-sided
$p=\mktLWp$. The univariate spanning alpha is \mktUniA\% a year
($t=\mktUniT$, $\beta=\mktUniB$).

The gain is not an alpha with respect to the factors. Over \mktFullSpan{},
regressing the managed market on all seven factors gives $\alpha=\mktFullA$\%
($t=\mktFullT$, $R^2=\mktFullRsq$). The loadings are \mktFullBMkt{} on
the market, \mktFullBRMW{} on RMW, \mktFullBMom{} on momentum, \mktFullBBAB{}
on BAB, \mktFullBCMA{} on CMA, \mktFullBSMB{} on SMB and \mktFullBHML{} on
HML. The Romano--Wolf adjusted $p$-value for the market is \mktRWp.

The loadings have a natural reading. The drawdown brake and the 21-day tilt
cut exposure after the market falls, which is a time-series trend component
and matches the momentum loading. De-levering when volatility is high is a
defensive tilt, and it co-moves with profitability and BAB. Once an investor
holds those factors, the overlay adds nothing.

\begin{table}[t]
\centering\small
\caption{Univariate alpha of the managed market by era (\% a year,
Newey--West $t$).}
\label{tab:eras}
\begin{tabular}{lrrr}
\toprule
Era & $\alpha$ & ($t$) & Months \\
\midrule
\tabEras
\bottomrule
\end{tabular}
\end{table}

Table~\ref{tab:eras} shows where the univariate alpha comes from. It is
positive in all \erasPos{} eras and concentrated before 1963: only
1927--1945 reaches $t\approx2$. In the two eras since 1963 it is about 1\% a
year with $t<1$. The sup-Wald statistic is \swStat, at \swAt, with $p=\swP$.

The contaminated window, \conSpan, is not part of the decision. It points the
same way: the univariate alpha is \conA\% ($t=\conT$), and the managed market
trails the market (\conSRm{} against \conSRr).

\paragraph{Grid, PBO and DSR.} Across the registered eight-member grid,
holdout Sharpe ratios range from \gridSRmin{} to \gridSRmax. The deployed
member is the highest, at \gridSRprim. The CSCV probability of backtest
overfitting across the grid is \gridPBO, over \gridCombos{} splits: the
in-sample winner usually lands below the out-of-sample median. The choice
among the grid's tilt, cap and brake settings therefore does not persist,
and the deployed member's rank is not evidence for those choices.

The deflated Sharpe ratio at $N=\nTrials$ is \dsrProg{}, with
SR$_0=\dsrSRzero$ taken from the grid's dispersion. It is uninformative
here: it tests whether the managed market's Sharpe ratio exceeds zero, and
the market's own Sharpe ratio does.

For the record, the platform's repository-wide PBO is \repoPBO{} over
\repoVariants{} variants.

\paragraph{Robustness.}
\begin{itemize}
\item \emph{Bill-rate financing.} Financing leverage at DTB3 instead of RF
  (from 1954, mean spread \dtbSpread\% a year) leaves the full-model alpha at
  \dtbFullA\% ($t=\dtbFullT$).
\item \emph{The volatility forecast.} The overlay's HAR forecast tracks VIX
  closely: the daily correlation is \vixHold{} in 1990--2012 and \vixCon{}
  after.
\end{itemize}

\subsection{Costs}

The overlay turns over \turnover{} units of exposure a year. Its mean
exposure is \meanExposure, and it sits at the cap \shareCap\% of the time.

At the univariate alpha, the breakeven one-way cost is \breakevenUni{}~bp per
unit of turnover. Against the full-model alpha the breakeven is negative:
no cost is low enough.

For comparison, the E-mini S\&P~500 has a minimum price fluctuation of 0.25
index points on a \$50-per-point contract (CME Rulebook, chapter~358). Half a
tick is $0.125/L$ of notional, about \emHalfTickBp{}~bp at an index level $L$
near 6{,}000.

Costs are therefore not what is missing on the market; the alpha is. Before
1982 no index future existed, and we do not model historical trading costs.

\subsection{Factor portfolios}

Table~\ref{tab:family} reports the family. These results are reported under
the registered Romano--Wolf correction and are not part of the decision.
\begin{itemize}
\item \emph{Survivors.} The full-model alphas of HML, RMW, CMA, momentum and
  BAB survive at 5\% family-wise. SMB's does not.
\item \emph{SPA.} Hansen's SPA test rejects the hypothesis that no managed
  factor has a positive alpha ($p=\spaP$).
\item \emph{Sharpe ratios.} They rise substantially for momentum and BAB,
  and for HML and RMW.
\end{itemize}

The rules are real-time: the scale is causal and there is no full-sample
constant. That distinguishes these results from the in-sample spanning
alphas. The pattern, with momentum and BAB strongest, matches the real-time
evidence of \citet{cederburg2020} and the momentum results of
\citet{barroso2015}.

\begin{table}[t]
\centering\small
\caption{The deployed overlay on factor portfolios, each first expressed in
the deployed volatility units. Alphas in \% a year, Newey--West $t$.
RW: Romano--Wolf stepdown-adjusted $p$ over the seven full-model alphas
(\rwSpan, $N=\rwN$). SR: Sharpe ratio of the raw and managed factor over
each factor's own span; LW $p$ one-sided. Scale: median $0.13/\hat\sigma$
before the overlay's own lever.}
\label{tab:family}
\begin{tabular}{lrrrrrrrrr}
\toprule
& \multicolumn{2}{c}{Univariate} & \multicolumn{2}{c}{Full model} & & \multicolumn{3}{c}{Sharpe ratio} & \\
\cmidrule(lr){2-3}\cmidrule(lr){4-5}\cmidrule(lr){7-9}
Factor & $\alpha$ & ($t$) & $\alpha$ & ($t$) & RW $p$ & Raw & Managed & LW $p$ & Scale \\
\midrule
\tabFamily
\bottomrule
\end{tabular}
\end{table}

Two caveats bound what Table~\ref{tab:family} supports.
\begin{enumerate}
\item \emph{Levered units.} The alphas are in levered units: the median
  scale runs from \scaleMin$\times$ to \scaleMax$\times$ before the
  overlay's own lever. The $t$-statistics and Sharpe ratios are scale-free;
  the alpha magnitudes are not.
\item \emph{No trading costs.} A long--short factor rebalances its stock
  legs every month and pays to short, and the overlay adds rescaling
  turnover on top. None of this can be priced from factor returns
  \citep[compare][]{novymarx2016,barroso2021}.
\end{enumerate}
We therefore report these as gross, real-time spanning alphas, and make no
net-of-cost claim.

\section{What the missing data would change}

The study used only data an undergraduate can obtain without a WRDS
subscription. Three missing datasets would sharpen specific results, and in
each case we stop rather than substitute.
\begin{itemize}
\item \emph{Stock-level CRSP} would let the overlay run on a point-in-time
  version of the deployed equal-weight book itself, rather than on the
  market as a proxy for it. It would also let the factor legs be rebuilt
  with their turnover, so that the family's alphas could be costed.
\item \emph{TAQ} would supply consolidated effective spreads against the
  NBBO, for SPY and for the factor legs. For the market overlay this cannot
  change the verdict, because no cost rescues a negative full-model alpha.
  For the family it is the missing half of a net-of-cost claim.
\item \emph{IEX HIST} trade messages would give a one-venue SPY spread.
  Measuring it requires accepting IEX's data terms, which we have not yet
  done. As above, it could not change the market verdict.
\end{itemize}

Two external checks are pending:
\begin{itemize}
\item the correlation of the overlay with S\&P's published daily
  risk-control index levels;
\item a broker paper account trading the overlay on SPY, which gives every
  order a third-party timestamp.
\end{itemize}

\section{Conclusion}

A capped, braked, forecast-driven volatility-targeting overlay, run
unchanged on a century of survivorship-free data it was never tuned on,
raises the market's Sharpe ratio. It does not earn alpha with respect to
the factors an investor can already hold. On the market it is risk control.

On factor portfolios the same real-time rule produces gross alphas that
survive multiple-testing correction. Whether they survive trading costs is a
question for stock-level data.

\appendix
\section{Reproducibility}

The registration, both amendments, the code, the results files and the
verification gates are in the public repository. The pre-registration is
\texttt{docs/design/DESIGN29.md}, and its amendments record the sample
correction and the seal.

The data manifest \texttt{data/vendor/MANIFEST.json} lists each vendor file's
URL, retrieval time, size and SHA-256.

Running \texttt{run\_research.py volalpha\_replication} and
\texttt{run\_research.py volalpha} rebuilds both results files from the
archived bytes. Two gates check the work:
\begin{itemize}
\item \emph{X46} checks the replication machinery and applies the
  registered tolerances.
\item \emph{X47} checks the overlay identity and the size and power of every
  test. It also refuses a results file not produced by the sealed code.
\end{itemize}
Every number in this paper is a macro generated from those files.

\bibliographystyle{plainnat}
\bibliography{volalpha}

\end{document}
